% Käsitsi ümber kirjutatud kogumiku lahendusest L11 (lk 37). Automaatne
% import märkis selle käsitsi tehtavaks hindamisskeemi tõttu; tekst ise on
% lühike proosa. Originaalis on teljestik kirjutatud kujul "mT", mis on
% m--T teljestik (mass horisontaalteljel, periood vertikaalteljel).
Statiivile kinnitatud vedrule kinnitatakse koormis ja leitakse stoppkella
abil $n$ võnkeks kulunud aeg $t$, millest saadakse võnkeperiood
$T = t/n$. Katset korratakse koormise erinevate massidega. Pendli
võnkeperioodi sõltuvus koormise massist esitatakse katseandmete põhjal
tabelina ning graafiliselt millimeetripaberil, teljestikus $m$--$T$.

\textbf{Hindamisskeem} (kokku 10 p): katse kirjeldus ja koormiste valik
\pp{1}; võnkeperioodi määramine seosest $T = t/n$ \pp{2};
kordusmõõtmiste kasutamine \pp{2}; tabeli vormistamine ja andmete kandmine
tabelisse \pp{1}; kordusmõõtmistest aritmeetilise keskmise arvutamine
\pp{1}; graafiku telgede õige määramine \pp{1}; graafiku mõõtkava \pp{1};
graafiku joonistamine \pp{1}.
